\documentclass[11pt]{article}

\usepackage[margin=1in]{geometry}
\usepackage{amsmath,amssymb,amsthm}
\usepackage{booktabs}
\usepackage[hidelinks]{hyperref}
\urlstyle{same}
\emergencystretch=2em

\newtheorem{theorem}{Theorem}[section]
\newtheorem{proposition}[theorem]{Proposition}
\newtheorem{lemma}[theorem]{Lemma}
\newtheorem{corollary}[theorem]{Corollary}
\newtheorem*{thmA}{Theorem A}
\theoremstyle{definition}
\newtheorem{definition}[theorem]{Definition}
\theoremstyle{remark}
\newtheorem{remark}[theorem]{Remark}

\newcommand{\N}{\mathbb N}
\newcommand{\F}{\mathbb F}
\newcommand{\Pow}{\operatorname{Pow}}
\newcommand{\opts}{\operatorname{opts}}
\newcommand{\cl}{\operatorname{cl}}
\newcommand{\xor}{\mathbin{\mathsf{xor}}}

\title{Natural Mis\`ere Quotients:\\
Exact Octal Certificates, Heap Normal Forms, and Grundy's Game}
\author{GPT 5.6 Sol \and a9lim}
\date{}

\begin{document}
\maketitle

\begin{abstract}
Siegel's valid-transition-table theorem characterizes abstract finite
mis\`ere quotients.  Natural realization asks for more: one uniform heap rule
must generate the table at every heap size.  We express this homogeneity as
an exact quotient-valued trace recurrence.  A least-rank separation argument
proves that the option-value set determines its value in every reduced valid
table.  Every value in such a table also has a context reaching the
$P$-portion, so no transition value can occur among its own options.  Hence a
fixed no-split code with finite quotient has an ultimately periodic pretending
function and a finite exact verification procedure.  Splitting introduces an
unbounded convolution, but an asserted periodic split trace still has a finite
certificate.  We identify the convolution exactly with a system of unary
Boolean language equations.  Without the valid-table axioms, totality is
$\Pi^0_1$-complete already for the single fixed code $0.7$; thus finiteness and
commutativity alone cannot yield an automaton argument.

We also prove a comparison theorem: every nontrivial finite quotient has an
exact one-species heap realization with finitely many source-local exceptions
and the uniform tail move $H_n\to H_{n-1}$.  This is not an octal realization;
indeed its inert pad cannot be retained in a verbatim octal encoding once the
quotient has more than two elements.  For every
$n\ge2$, the code
with $2^{n-1}$ consecutive digits equal to $3$ has quotient $\mathcal T_n$,
the tame quotient of order $2^n+2$; the proof gives the quotient map and the
outcome strategy for arbitrary heap multiplicities.  For mis\`ere Grundy's
game, the quotient generated by heaps through $12$ is $\mathcal T_2$, heap
$13$ enlarges it to a reduced monoid of order $12$, and heap $18$ enlarges it
to one of order $24$.  These are exact submonoid theorems, not bounded
indistinguishability computations.  If the full Grundy quotient is finite,
its kernel is necessarily nontrivial.  The full quotient and the uniform
realization of the wild $|P|=2$ family remain open.
\end{abstract}

\section{Introduction}

An abstract finite impartial game can be assembled to realize a prescribed
transition table.  A finite octal code cannot be assembled heap by heap: its
$k$th digit imposes the same translation-invariant move on every sufficiently
large heap.  The natural realization problem is to identify the extra
algebraic condition imposed by that uniform tail.

The proof architecture has four layers.  Reduced transition tables are
deterministic in their option-value sets and forbid self-options.  The octal
rule turns those facts into an exact recurrence and obstruction calculus for
the single-heap quotient word.  A finite-exception normal form separates
abstract heap realization from strict octal homogeneity.  Finally, the trace
recurrence can be solved either uniformly, as for the tame family, or prefix
by prefix with all multiplicities retained, as for Grundy's game.  The main
results are the following.

\begin{thmA}
\begin{enumerate}
\item A reduced valid transition table assigns at most one value to each
option-value set.  A finite octal code realizes a finite reduced bipartite
monoid exactly when its induced heap trace is total, lies in such a table,
and generates the monoid.
\item No table transition has its value among its option values.  Thus every
whole-, one-, and two-remainder octal digit imposes an explicit inequality on
the heap trace.
\item If the code has no split bits and its quotient is finite, the heap trace
is ultimately periodic.  Exact realization by a fixed no-split code is
therefore decidable.  If a split trace is already ultimately periodic, its
infinite tail is likewise certified by finitely many records.
\item For an arbitrary finite partial decoder, fixed-code split totality is a
unary Boolean-language-equation problem and is $\Pi^0_1$-complete already for
$0.7$.  The hardness construction is not a valid mis\`ere transition table;
validity is the exact remaining separator.
\item Every nontrivial finite quotient has an exact finite-exception unary-heap
normal form: finitely many source-local prefix rules followed by the uniform
tail $H_n\to H_{n-1}$.  For quotients of order greater than two, its inert pad
rules out a verbatim finite-octal encoding of the chosen prefix records.
\item For every $n\ge2$, the code
\(0.\underbrace{33\cdots3}_{2^{n-1}\text{ digits}}\) has exact quotient
$\mathcal T_n$.  Thus every tame quotient in the finite $|P|=2$
classification has a finite-code realization.
\item In mis\`ere Grundy's game the exact quotients generated through heaps
$12$, $13$, and $18$ have orders $6$, $12$, and $24$, respectively, with
the presentations and reduced translation tables given below.  Any finite
full Grundy quotient has a nontrivial kernel.
\end{enumerate}
\end{thmA}

The first four items give exact algebraic, finite-certificate, and
computability boundaries; the fifth is a realization theorem in a comparison
class, not for finite octal codes; the sixth is a uniform octal realization
theorem; and the seventh combines exact prefix theorems with a conditional
structure theorem.  None asserts that valid split traces are automatically
periodic or that the full Grundy quotient is finite.
Sections~\ref{sec:tables} and~\ref{sec:traces} build the trace calculus,
Section~\ref{sec:normalform} gives the comparison normal form,
Section~\ref{sec:tame} realizes the tame family, and
Section~\ref{sec:grundy} gives the Grundy extensions.
Section~\ref{sec:boundary} states the formal and open boundaries.

\section{Transition tables and deterministic values}\label{sec:tables}

Let $A$ be a set of finite impartial games closed under options and
disjunctive sum.  Under mis\`ere play the terminal game is an $N$-position.
For $G,H\in A$, put
\[
 G\sim_A H
 \quad\Longleftrightarrow\quad
 o^-(G+X)=o^-(H+X)\quad\hbox{for every }X\in A.
\]
The quotient $Q(A)=A/{\sim_A}$ is a commutative monoid.  It comes with the
distinguished subset $P$ of previous-player wins.  The pair $(Q,P)$ is
\emph{reduced}: if $x\ne y$, then some $z$ has exactly one of $xz,yz$ in
$P$.  Also $1\notin P$, because the terminal position is $N$.

Plambeck and Siegel developed this quotient method and the generalized mex
rule in \cite{plambecksiegel2008}; Siegel subsequently proved the abstract
classification by valid transition tables in \cite{siegel2007}.  A
transition pair is $(x,E)\in Q\times\Pow(Q)$, with product
\begin{equation}
 (x,E)(y,F)=(xy,yE\cup xF).
\end{equation}
A transition table $T$ on $(Q,P)$ is valid when it is complete, closed under
the displayed product, parity-correct,
\begin{equation}
 (x,E)\in T
 \quad\Longrightarrow\quad
 x\in P\ \Longleftrightarrow\ E\ne\varnothing\text{ and }E\cap P=\varnothing,
\end{equation}
and well-founded: there is $R:Q\to\N$, with $R(1)=0$, such that every
$x$ has some $(x,E)\in T$ satisfying $R(e)<R(x)$ for all $e\in E$.
Siegel's theorem says that a reduced $(Q,P)$ with $1\notin P$ is an abstract
finite mis\`ere quotient if and only if it admits such a table.

The natural problem begins after that theorem: a code must generate its table
through one translation-invariant heap rule.

\subsection{Option-set determinism}

The first observation converts the natural problem from a nondeterministic
labeling problem into a trace recurrence.

\begin{theorem}[Option-set determinism]\label{thm:determinism}
Let $T$ be a valid transition table on a reduced bipartite monoid $(Q,P)$.
If
\[
 (x,E)\in T\quad\text{and}\quad(y,E)\in T,
\]
then $x=y$.
\end{theorem}

\begin{proof}
Assume $x\ne y$.  Among all contexts separating $x$ and $y$, choose $z$
with $R(z)$ least.  Take a descending pair $(z,F)\in T$.  Suppose, after
interchanging $x,y$ if necessary, that $xz\in P$ and $yz\notin P$.
Closure gives
\[
 (xz,zE\cup xF),\qquad (yz,zE\cup yF)\in T.
\]
Parity for $xz$ says that $zE\cup xF$ is nonempty and contains no element
of $P$.  The second option set is nonempty as well: a witness coming from
$zE$ is unchanged, while a witness $xf$ is replaced by $yf$.
Since $yz\notin P$, parity therefore gives some $P$-element in
$zE\cup yF$.  It cannot lie in $zE$, which is already contained in the
$P$-free first option set.  Hence $yf\in P$ for some $f\in F$, while
$xf\notin P$.  Thus $f$ also separates $x,y$, but $R(f)<R(z)$, a
contradiction.
\end{proof}

Equivalently, every valid table defines a partial function
\begin{equation}
 \mu_T(E)=x\quad\Longleftrightarrow\quad (x,E)\in T.
\end{equation}
This strengthening of the option-image principle is purely algebraic: it
does not assume that $T$ was first obtained from games.  The theorem and its
least-rank proof are kernel-checked in the Lean module
\texttt{Ogdoad.MisereTransition}.

We also record the standard local obstruction in a form suited to heap
traces.  Define the meximal set
\[
 M_x=\{y\in Q:\text{there is no }z\text{ with }xz,yz\in P\}.
\]

\begin{lemma}[Meximal obstruction]\label{lem:meximal}
Every $(x,E)\in T$ satisfies $E\subseteq M_x$.
\end{lemma}

\begin{proof}
If $e\in E$ and $xz,ez\in P$, choose any table pair $(z,F)$.  Closure puts
$(xz,zE\cup xF)$ in $T$, and $ez\in zE$ is a $P$-valued option of the
$P$-value $xz$, contradicting parity.
\end{proof}

The fact needed to exclude self-options is stronger than reducedness alone;
validity supplies it.

\begin{proposition}[$P$-reaching contexts and loop-free records]
\label{prop:noself}
Let $T$ be a valid transition table on a nontrivial reduced bipartite monoid
$(Q,P)$.  For every $x\in Q$ there is a $z\in Q$ with $xz\in P$.
Consequently every $(x,E)\in T$ satisfies $x\notin E$.
\end{proposition}

\begin{proof}
Choose a descending record $(x,E)$.  If $E=\varnothing$, determinism against
the terminal record $(1,\varnothing)$ gives $x=1$.  Nontrivial reducedness
implies $P\ne\varnothing$, so any $p\in P$ is a $P$-reaching context.
Suppose $E\ne\varnothing$.  If $x\in P$, take $z=1$; if $x^2\in P$, take
$z=x$.  Otherwise closure puts $(x^2,xE)$ in $T$.  Its option set is
nonempty, so parity and $x^2\notin P$ force $xe\in P$ for some $e\in E$.
Thus every value has a $P$-reaching context.

The meximal obstruction gives $E\subseteq M_x$.  If $x\in E$, then
$x\in M_x$, contradicting a context $z$ for which both copies of $xz$ lie
in $P$.
\end{proof}

The one-element quotient has only the terminal record and satisfies the same
no-self conclusion trivially.  Reducedness without a valid table would not be
enough for the $P$-reaching statement.

\section{Exact octal traces}\label{sec:traces}

\subsection{The trace criterion}

Write a finite octal code as three masks.  Removing $k$ counters may take a
whole heap when $C_k$ holds, may leave one nonempty heap when $A_k$ holds,
and may leave two nonempty heaps when $B_k$ holds.  These are respectively
the $1,2,4$ bits of the $k$th octal digit.  Let $d$ be the last nonzero
digit, put $x_0=1$, and seek $x_n=\Phi(H_n)$ for $n\ge1$.  The quotient
image of the option set of a heap is exactly
\begin{align}
 E_n={}&\bigl(\{1\}\text{ if }C_n\bigr)
 \cup\{x_{n-k}:A_k,\ k<n\}\\
 &\cup\{x_i x_{n-k-i}:B_k,\ k\le n-2,\
                 1\le i\le (n-k)/2\}.\nonumber
\end{align}
Equal parts are permitted in an octal split.

Proposition~\ref{prop:noself} immediately gives three absolute exclusions for
every exact trace:
\begin{align}
 C_n&\Longrightarrow x_n\ne1,\\
 A_k,\ k<n&\Longrightarrow x_n\ne x_{n-k},\\
 B_k,\ i,j\ge1,\ i+j=n-k
   &\Longrightarrow x_n\ne x_i x_j.
\end{align}
These are not bounded-search heuristics: the indicated value occurs in
$E_n$, so equality would make the heap record a forbidden self-option.  For
example, if the tail has period $p$, no active one-remainder offset can be a
multiple of $p$; applying the second exclusion sufficiently far into the tail
gives an immediate contradiction.

\begin{theorem}[Exact trace criterion]\label{thm:trace}
Fix a finite octal code $\Gamma$ and a finite reduced bipartite monoid
$(Q,P)$.  Then $Q(\Gamma)\cong(Q,P)$ if and only if there are a valid table
$T$ on $(Q,P)$ and a sequence $(x_n)_{n\ge0}$ such that
\begin{enumerate}
\item $x_0=1$;
\item $(x_n,E_n)\in T$ for every $n\ge1$, with $E_n$ given by
      the displayed octal-trace formula;
\item the elements $x_n$ generate $Q$.
\end{enumerate}
For fixed $T$ and $\Gamma$, the sequence is unique whenever it exists.
\end{theorem}

\begin{proof}
Necessity follows by taking the actual transition algebra and the quotient
images of the heaps.  Conversely define $\Phi$ multiplicatively on heap
multisets.  The transition pair of any position is the product of its heap
pairs.  Closure of $T$ therefore puts
$(\Phi(G),\Phi''G)$ in $T$ for every position $G$.  Induction on total
counter count, using parity, gives
\[
 o^-(G)=P\quad\Longleftrightarrow\quad\Phi(G)\in P.
\]
Generation makes $\Phi$ onto.  Equal images are indistinguishable by the
displayed equivalence; distinct images are distinguished by a context in
$Q$, and surjectivity realizes that context by a position.  Thus the fibers
are exactly the indistinguishability classes.  Uniqueness is
Theorem~\ref{thm:determinism} applied recursively.
\end{proof}

This theorem identifies the general unresolved core precisely: a finite
alphabet quadratic-convolution recurrence must remain defined forever.  It
also separates exact certificates from bounded signatures.

\subsection{Finite-state and periodic traces}

Suppose every digit lies in $\{0,1,2,3\}$, so $B_k$ is always false.  For
$n>d$, whole-heap removal is also inactive, and
\begin{equation}
 E_n=\{x_{n-k}:A_k,\ 1\le k\le d\}.
\end{equation}

\begin{theorem}[Automatic periodicity]\label{thm:periodicity}
Let $\Gamma$ be a finite no-split octal code with last nonzero digit $d$.
If $Q(\Gamma)$ is finite, then its single-heap quotient word
$x_n=\Phi(H_n)$ is ultimately periodic.  A repeated state occurs among at
most $|Q|^d+1$ consecutive $d$-windows after the exceptional prefix.
\end{theorem}

\begin{proof}
The ordered window
\[
 W_n=(x_{n-d},\ldots,x_{n-1})\in Q^d
\]
determines $E_n$ by the preceding formula.  By
Theorem~\ref{thm:determinism}, $E_n$ determines $x_n$.  Thus $W_n$
determines $W_{n+1}$.  The deterministic map acts on at most $|Q|^d$
states, so one state repeats and the entire subsequent orbit is periodic.
\end{proof}

\begin{corollary}\label{cor:decidable}
For a fixed finite no-split code and a fixed finite reduced $(Q,P)$, exact
realization is decidable.  The same is true when the code length is bounded
in advance.
\end{corollary}

Indeed there are finitely many transition tables on $Q$; validity is a
finite condition, and each induced trace either fails or enters a cycle.
Generation and parity can then be checked on finite data.  The existential
problem over codes of unbounded length is not settled by this argument.

Splitting is the exact place where finite memory fails.  The new term
\[
 S_m=\{x_i x_j:i,j\ge1,\ i+j=m\}
\]
depends on an unbounded convolution of the prefix.  Nevertheless an already
periodic trace remains finitely certifiable.

\begin{proposition}[Periodic convolution]\label{prop:convolution}
Assume $x_{n+p}=x_n$ for all $n\ge N$, where $N,p\ge1$.  Then
\[
 S_{m+p}=S_m\qquad(m\ge2N+p).
\]
Consequently the records $(x_n,E_n)$ of a length-$d$ octal code are
periodic for all $n\ge2N+p+d$.
\end{proposition}

\begin{proof}
For a pair $i+j=m+p$, one of $i,j$ is at least $N+p$; subtract $p$ from
that component and use periodicity.  Conversely, for $i+j=m$, one component
is at least $N$; add $p$ to it.  Commutativity permits reordering.  The
one-remainder terms are easier, and whole-heap bits are inactive past $d$.
\end{proof}

More explicitly, if the code is supported through digit $d$, then checking
one full period of records after any bound $B\ge2N+p+d$ certifies every later
heap record.  Thus a proposed preperiod, period, valid table, finite trace
through that first tail period, and generation check form an exact
infinite-rule certificate.  This is a transition-table form of the
periodicity argument in \cite{siegel2006notes,plambecksiegel2008}.

\subsection{The Boolean-language boundary}\label{sec:languages}

The convolution obstruction can be made exact.  For languages
$L,M\subseteq\N$, write
\[
 L+M=\{i+j:i\in L,\ j\in M\}.
\]
Fix a candidate word $x=(x_n)$ and put $X_q=\{n:x_n=q\}$.  For each
$q\in Q$, let $A_q$ be the language of heap sizes at which $q$ occurs as an
option value.  In terms of the three masks it is
\begin{equation}
 A_q=C_q\ \cup\!\bigcup_{k\in A}\bigl(k+(X_q\cap\N_{>0})\bigr)\ \cup\!
 \bigcup_{k\in B}\ \bigcup_{uv=q}
       \bigl(k+(X_u\cap\N_{>0})+(X_v\cap\N_{>0})\bigr),
\end{equation}
where $C_1=C$ is the finite whole-removal language and $C_q=\varnothing$
for $q\ne1$.  Bounds excluding empty remainders are understood in the
shifted terms.  For $S\subseteq Q$, define the exact Boolean cell
\begin{equation}
 D_S=\bigcap_{q\in S}A_q\ \cap\!
     \bigcap_{q\notin S}(\N\setminus A_q).
\end{equation}

\begin{proposition}[Exact language equations]\label{prop:language-equations}
Let $\mu:\Pow(Q)\rightharpoonup Q$ be any finite partial decoder.  A word
$x$ is its total octal trace exactly when
\begin{equation}
 X_q=\bigcup_{\mu(S)=q}D_S\qquad(q\in Q)
\end{equation}
and the failure language
\begin{equation}
 F_\mu=\bigcup_{\mu(S)\text{ undefined}}D_S
\end{equation}
is empty.  Coefficient $n$ belongs to $F_\mu$ exactly when the decoder has no
value on the complete option set at $n$.
\end{proposition}

\begin{proof}
The first displayed equation is simply the whole-, one-, and
two-remainder definition sorted by quotient value.  Hence $n\in D_S$ if and
only if the complete option-value set at $H_n$ is exactly $S$.  Applying
$\mu$ coefficientwise gives both displayed conclusions.
\end{proof}

These equations are unary, but they are not finite-state equations.  Unary
conjunctive grammars already generate nonregular languages such as
$\{a^{4^r}:r\ge0\}$ \cite{jez2008}.  In fact the unconstrained trace problem
reaches the full effective boundary.

\begin{theorem}[Arbitrary-decoder hardness]\label{thm:decoder-hardness}
For the single fixed octal code $0.7$, the following problem is
$\Pi^0_1$-complete: given a finite commutative monoid $M$ and an arbitrary
finite partial decoder $\mu:\Pow(M)\rightharpoonup M$, decide whether the trace
starting from $x_0=1$ is total.
\end{theorem}

\begin{proof}
Non-totality has a least failing coefficient, so it is recursively
enumerable and totality is in $\Pi^0_1$.

For hardness, start with a unary unambiguous conjunctive grammar.  Emptiness
for these grammars is undecidable by an effective reduction from Turing
machine emptiness \cite[Lemma~12]{jezokhotin2017}.  Effectively put the
grammar in epsilon-free binary conjunctive normal form and prefix a fixed
two-letter word, so that lengths $0$ and $1$ are irrelevant while emptiness
is preserved.  Unambiguity is needed only for the cited source theorem; the
normalization need not preserve it.  Let $\mathcal N$ be the resulting finite
nonterminal set and let $V_n\subseteq\mathcal N$ be its truth vector on
$a^n$.  Take the full, effectively known Boolean cube
$C=\Pow(\mathcal N)$ as the color set and form the commutative monoid
\[
 M_C=\{\text{multisets on $C$ of degree at most $2$}\}\sqcup\{\bot\}.
\]
Multiplication is multiset union until degree $2$ and is $\bot$ after
overflow, with $\bot$ absorbing; the empty multiset is the identity.  This is
associative: for three ordinary factors both parenthesizations give their
union when its total degree is at most two and $\bot$ otherwise, while an
existing $\bot$ remains absorbing.

Assign $H_h$ the singleton marker $[V_{h+1}]$.  Under $0.7$, the option set
of $H_h$ for $h\ge3$ is exactly
\begin{equation}
 \{[V_h]\}\ \cup
 \{[V_{i+1},V_{j+1}]:i,j\ge1,\ i+j=h-1\}.
\end{equation}
The degree-two markers encode every binary cut of the word of length $h+1$
whose two parts have length at least two.  The degree-one predecessor marker,
together with the fixed vector $V_1$, encodes the two cuts having a part of
length one.  Because the alphabet is unary, unordered pairs lose no cut
information.  Thus the exact option set determines $V_{h+1}$.

More explicitly, the normalized grammar gives an effective finite update
map
\[
 \mathsf H:C\times\Pow(\{[u,v]:u,v\in C\})\longrightarrow C.
\]
It computes the next truth vector from the predecessor vector and the set of
unordered binary-cut vectors; the two cuts involving $V_1$ are folded into
the update.  For every decodable signature
\[
 S=\{[c]\}\cup R,
 \qquad R\subseteq\{[u,v]:u,v\in C\},
\]
define $\mu(S)=[\mathsf H(c,R)]$ when the start nonterminal is absent from
$\mathsf H(c,R)$, and leave $\mu(S)$ undefined when it is present.  This is a
uniform finite definition, independent of which signatures the eventual
trace reaches.  Malformed signatures may be left undefined.  The genuinely
exceptional $H_1$ cell computes $V_2$ directly.  The $H_2$ cell is the
decodable signature $(c,R)=(V_2,\varnothing)$, for which
$\mathsf H(V_2,\varnothing)=V_3$ because only the two length-one cuts occur.
Both cells are left undefined exactly when their output start bit is true.

Before the first accepted length, induction now gives precisely the asserted
markers.  At the first accepted length the decoder fails.  Hence the trace is
total if and only if the grammar language is empty.  Turing-machine emptiness
is $\Pi^0_1$-complete, completing the reduction.
\end{proof}

The decoder in Theorem~\ref{thm:decoder-hardness} is deliberately arbitrary.
It need not extend to a reduced, closed, parity-correct, ranked mis\`ere
transition table.  Therefore the theorem neither proves undecidability of
natural realization nor decides a fixed valid-table trace.  It proves the
more surgical statement needed here: any decision argument must use the
valid-table axioms in a load-bearing way.  Finiteness of $Q$, commutativity of
its multiplication, and unary input alone are insufficient.

\section{A finite-exception heap normal form}\label{sec:normalform}

The abstract table theorem does admit a uniform one-species heap realization
once the finite prefix is allowed to be source-local.  This comparison result
pinpoints why the direct construction does not immediately become octal.

\begin{theorem}[Finite-exception unary-heap normal form]
\label{thm:normalform}
Let $(Q,P)$ be a nontrivial finite reduced bipartite commutative monoid.
Suppose $T$ is a closed parity-correct transition table and $R:Q\to\N$, with
$R(1)=0$, such that every $q\in Q$ has a record $(q,E_q)\in T$ satisfying
$R(e)<R(q)$ for all $e\in E_q$.  Then there are $N$ and a finite-description
impartial ruleset on numerical heaps $H_1,H_2,\ldots$ such that
\begin{enumerate}
\item for $n\le N$, every option of $H_n$ is a single heap $H_j$ with $j<n$;
\item $H_N$ is terminal;
\item for every $n>N$, $\opts(H_n)=\{H_{n-1}\}$; and
\item the exact mis\`ere quotient of all finite sums of these heaps is
      isomorphic to $(Q,P)$.
\end{enumerate}
\end{theorem}

\begin{proof}
First find a star record.  Choose a nonidentity value $a$ of least rank and a
descending record $(a,E)$.  Every member of $E$ has smaller rank, hence is
$1$.  The set $E$ cannot be empty, since determinism against
$(1,\varnothing)$ would give $a=1$.  Therefore $(a,\{1\})\in T$.

For each $q\in Q$, choose a descending record $(q,E_q)$.  Enumerate the
elements of $Q$ in nondecreasing rank and make one prefix heap $B_q$ with
options $B_e$ for $e\in E_q$.  Every target is an earlier heap.  Append one
inert padding heap $I$, followed by heaps $U_0,U_1,\ldots$ with
\[
 \opts(U_0)=\{I\},\qquad \opts(U_{k+1})=\{U_k\}.
\]
Assign the values
\[
 \Phi(B_q)=q,\qquad \Phi(I)=1,\qquad
 \Phi(U_k)=a^{k+1}.
\]
Closure of $T$ and $(a,\{1\})\in T$ give by induction
\[
 (a^{k+1},\{a^k\})\in T,
\]
so every single-heap record lies in $T$.  The transition record of a sum is
the product of its component records and therefore also lies in $T$.
All moves strictly lower the sum of the displayed heap ranks.  Induction on
that rank and parity show
\[
 o^-(G)=P\quad\Longleftrightarrow\quad\Phi(G)\in P
\]
for every position $G$.  Each $q$ is represented by $B_q$, so $\Phi$ is
surjective.  Equal values are indistinguishable, while reducedness and
surjectivity realize a separating context for unequal values.  Hence the
fibers of $\Phi$ are exactly the indistinguishability classes.
\end{proof}

This is an eventually translation-invariant finite-exception heap normal
form, not a finite-octal realization.  Turning a prefix edge
$H_i\to H_{i-d}$ into an octal digit at offset $d$ also creates translated
moves from every larger heap.  The construction supplies no values for those
additional targets.  Controlling that cross-talk---possibly by a different,
generator-level encoding---remains an unresolved bridge.  The comparison
class and finite octal games are therefore not ordered by inclusion: the
former permits arbitrary source-local prefix moves but only a unary tail,
whereas the latter permits persistent multi-removal and split moves.

The inert pad is not merely an expositional nuisance.

\begin{proposition}[Inert-pad obstruction]\label{prop:pad-obstruction}
Assume $|Q|>2$, and choose one strictly descending record $(q,E_q)$ for each
$q\in Q$ as in Theorem~\ref{thm:normalform}.  There is no finite octal code
with prefix heaps $B_q$ followed by a later terminal heap $I$ such that every
$B_q$ has exact option-value set $E_q$.
\end{proposition}

\begin{proof}
Determinism identifies every chosen empty record with
$(1,\varnothing)$ and identifies all chosen records with option set $\{1\}$
with the unique star value.  Since $|Q|>2$, some chosen $E_q$ therefore
contains an $e\ne1$.

A whole-heap-removal bit contributes only $1$, so the option $e$ at $B_q$
must come from a one-remainder or two-remainder bit at some removal size $k$.
Such a bit is persistent.  If it leaves one remainder at the earlier heap
$B_q$, then $k<|B_q|<|I|$ and it leaves a remainder at $I$.  If it leaves two
remainders at $B_q$, then $k+2\le |B_q|<|I|$ and it also gives a split option
at $I$.  Either conclusion contradicts that $I$ is terminal.
\end{proof}

Thus no placement trick, including spacing the representatives farther
apart, turns the proof of Theorem~\ref{thm:normalform} verbatim into an octal
proof.  This does not rule out octal universality: a successful compiler may
replace the inert pad by an active bridge or may encode quotient generators
contextually rather than representing every chosen descending record by one
heap.

\section{Uniform realization of the tame family}\label{sec:tame}

For $n\ge2$, let
\[
 m=2^{n-1},\qquad L=m+1,
\]
and let $\Gamma_n$ be the finite octal code
\[
 \Gamma_n=0.\underbrace{33\cdots3}_{m\text{ digits}}.
\]
A move subtracts any number in $\{1,\ldots,m\}$ from one heap, deleting it
if the remainder is zero and otherwise leaving the one remainder heap.
There are no splits.

Define
\[
 Q_n=\{1,a\}\sqcup\{z_v:v\in\F_2^n\},\qquad
 P_n=\{a,z_0\}.
\]
Identify $v\in\F_2^n$ with an integer in $[0,2^n)$ and write vector addition
as XOR.  Multiplication is
\begin{equation}
 a^2=1,\qquad z_u z_v=z_{u\xor v},\qquad
 a z_v=z_{v\xor1}.
\end{equation}
This is Siegel's tame quotient $\mathcal T_n$: its kernel is
$\{z_v\}\cong(\mathbb Z/2)^n$, and $|Q_n|=2^n+2$
\cite{plambecksiegel2008,siegel2006notes}.

\begin{theorem}[Tame-family realization]\label{thm:tame}
For every $n\ge2$, the exact mis\`ere quotient of $\Gamma_n$ is
$\mathcal T_n$.
\end{theorem}

\begin{proof}
For a heap $H_h$, let $r(h)$ be the residue of $h$ modulo $L$ and define
\begin{equation}
 \Phi(H_h)=
 \begin{cases}
  1,&r(h)=0,\\
  a,&r(h)=1,\\
  z_{r(h)},&2\le r(h)\le m.
 \end{cases}
\end{equation}
Extend $\Phi$ multiplicatively.

We first prove the outcome criterion directly.  Discard residue-zero heaps.
Let $t\in\F_2$ be the parity of the number of residue-one heaps, and let
$R$ be the multiset of residues in $\{2,\ldots,m\}$.  From
the displayed multiplication laws,
\begin{equation}
 \Phi(X)\in P_n\quad\Longleftrightarrow\quad
 \begin{cases}
  t=1,&R=\varnothing,\\
  t\xor\displaystyle\bigoplus_{r\in R}r=0,&R\ne\varnothing.
 \end{cases}
\end{equation}

Induct on total counter count.  A legal move changes one residue $r$ to a
different residue $s$.  If $R=\varnothing$ and $t=1$, changing a residue
$0$ or $1$ either flips the parity or creates a nonzero XOR, so every option
fails this criterion.  If $R\ne\varnothing$ and the displayed
XOR is zero, there must be at least two residues at least $2$: one such
residue cannot cancel either $0$ or $1$.  A move changes the XOR by
$r\xor s\ne0$, and at least one large residue remains.  Hence every option
of a predicted $P$-position is predicted $N$.

Conversely suppose the position is predicted $N$.  If $R=\varnothing$ and
$t=0$, move a residue-one heap to residue zero when one exists.  Otherwise
the nonempty position has a positive residue-zero heap; subtracting $m$
moves it to residue one.  Either move makes $t=1$.

Now suppose at least two residues at least $2$ occur, and write
\[
 w=t\xor\bigoplus_{r\in R}r\ne0.
\]
Choose a component residue $r$ whose binary expansion contains the highest
set bit of $w$ and put $s=r\xor w$.  Such a component exists among the
residues contributing to the XOR: it may be a residue-one component when
the highest bit is the low bit, and otherwise it lies in $R$.  Then $s<r$,
so subtracting
$r-s\le m$ is legal, the new XOR is zero, and at least one large residue
remains.  Finally, if exactly one large residue $r$ occurs, move it to
residue $1-t<r$.  No large residue remains and the new residue-one parity is
odd.  This proves the displayed criterion is the actual mis\`ere outcome
for every position.

The map is onto.  The heap residues provide $a$ and
$z_2,z_4,\ldots,z_m$; also $z_2^2=z_0$ and $az_0=z_1$.  These are a binary
basis for every $z_v$.  Finally $(Q_n,P_n)$ is reduced.  Distinct
$z_u,z_v$ are separated by $z_u$.  The elements $1,a$ are separated by
$1$; the context $a$ separates $1$ from $z_v$ unless $v=1$, when $z_1$
does; and the context $1$ separates $a$ from $z_v$ unless $v=0$, when
$z_0$ does.  Outcome correctness, surjectivity, and reduction identify the
fibers of $\Phi$ with exact indistinguishability classes.
\end{proof}

For $n=2$, the construction is $0.33$ and gives the six-element quotient
$\mathcal T_2$.  The theorem realizes the entire tame arm in Siegel's
classification of finite quotients with $|P|=2$; the other arm consists of
the tame extensions of the eight-element wild quotient $\mathcal R_8$
\cite{siegel2007}.  The base $\mathcal R_8$ itself is realized by $0.75$
\cite{plambecksiegel2008}, but a uniform finite-code realization of all its
tame extensions is not supplied here.

\section{Exact initial quotients of mis\`ere Grundy's game}\label{sec:grundy}

In Grundy's game a heap may be split into two unequal positive heaps:
\[
 \opts(H_n)=\{H_i+H_{n-i}:1\le i<n-i\}.
\]
The rule preserves total counters but decreases the well-founded quantity
$\sum_h\max(h-2,0)$.  Its mis\`ere form is a longstanding difficult example;
the early analysis goes back to Grundy and Smith \cite{grundysmith1956} and
is discussed in \cite{berlekampconwayguy2001}.

\subsection{The tame prefix through heap 12}

We first give an exact transition certificate through heap $12$.  Write
\[
 \mathcal T_2=\langle a,b\mid a^2=1,\ b^3=b\rangle,
 \qquad P=\{a,b^2\}.
\]
The heap values are
\[
 x_1=x_2=1,\qquad
 x_n=
 \begin{cases}
 a,&n\equiv0\pmod3,\\
 1,&n\equiv1\pmod3,\\
 b,&n\equiv2\pmod3
 \end{cases}
 \quad(3\le n\le12).
\]
Their distinct transition records are
\begin{align*}
 &(1,\varnothing),\ (a,\{1\}),\ (1,\{a\}),\
 (b,\{1,a\}),\\
 &(a,\{1,b\}),\ (1,\{a,b\}),\
 (b,\{1,a,ab\}).
\end{align*}
Closing these records under the transition product gives the following
$24$ pairs; each row lists the option sets paired with its left-hand value.
\begin{center}
\small
\begin{tabular}{c@{\quad}l}
\toprule
$1$ & $\varnothing$; $\{a\}$; $\{a,b\}$; $\{a,ab\}$;
       $\{a,b,ab\}$\\
$b$ & $\{1,a\}$; $\{b^2,ab^2\}$; $\{1,a,ab\}$;
       $\{b^2,ab,ab^2\}$; $\{1,b^2,a,ab\}$;
       $\{1,a,ab,ab^2\}$; $\{1,b^2,a,ab,ab^2\}$\\
$b^2$ & $\{b,ab\}$; $\{b,ab,ab^2\}$\\
$a$ & $\{1\}$; $\{1,b\}$; $\{1,ab\}$; $\{1,b,ab\}$\\
$ab$ & $\{1,b,a\}$; $\{b,b^2,ab^2\}$;
       $\{1,b,b^2,a\}$; $\{1,b,a,ab^2\}$;
       $\{1,b,b^2,a,ab^2\}$\\
$ab^2$ & $\{b,b^2,ab\}$\\
\bottomrule
\end{tabular}
\end{center}
The table is closed and every pair obeys parity; $a,b$ occur as heap values
and $\mathcal T_2$ is reduced.  The trace criterion therefore proves:

\begin{proposition}\label{prop:grundy12}
The exact mis\`ere quotient generated by $H_1,\ldots,H_{12}$ is
$\mathcal T_2$.
\end{proposition}

\subsection{The first extension at heap 13}

Before extending this quotient, we isolate a substitution point that is
easy to miss.  Write $G\equiv_-H$ when $G$ and $H$ have the same mis\`ere
outcome after addition of every finite impartial game, not merely every
game in the present heap universe.

\begin{lemma}[Global substitution]\label{lem:global-substitution}
Let $U$ be the additive closure of games $G_i$, and suppose that each $G_i$
is globally equivalent to a representative $R_i$ lying in an additively
closed subuniverse $V\subseteq U$.  Then inclusion induces an isomorphism
$Q(V)\cong Q(U)$.
\end{lemma}

\begin{proof}
Global equivalence is an additive congruence, so every $X\in U$ is globally
equivalent to the sum $r(X)\in V$ obtained by replacing its components.
If $A,B\in V$ agree against all $V$-contexts, then for every $X\in U$,
\[
 o^-(A+X)=o^-(A+r(X))=o^-(B+r(X))=o^-(B+X).
\]
The converse follows because $V\subseteq U$, and every $U$-class has a
representative in $V$.
\end{proof}

The Grundy--Smith canonical reduction rules are global reductions.  Applied
to the explicit option lists on heaps through $12$, they give
\[
 \begin{array}{c|l}
 0&H_1,H_2,H_4,H_7,H_{10}\\
 a=*1&H_3,H_6,H_9,H_{12}\\
 b=*2&H_5,H_8,H_{11}.
 \end{array}
\]
These are the reductions displayed in \cite[pp.~419--421]{berlekampconwayguy2001}.
Consequently adjoining $H_{13}$ cannot secretly refine one of these earlier
identifications: they remain valid against the new context $H_{13}$ and
all of its sums.

Heap $13$ has option-value set $\{a,b,b^2\}$.  The alleged continuation
$x_{13}=1$ is impossible: multiplying its transition by two copies of the
$H_5$ transition produces
\[
 (b^2,\{b,b^2,ab,ab^2\}),
\]
which violates parity because both the value and one option lie in $P$.
Concretely $H_4+2H_5$ is $P$, while $H_{13}+2H_5$ has the $P$-option
$H_8+3H_5$ and is $N$.

The extension can still be solved symbolically.  Let $c$ denote the new
heap value, with transition $c\to\{a,b,b^2\}$.  For raw monomials
$a^e b^j c^k$, direct induction on $(k,j,e)$ gives the following complete
outcome law:
\begin{equation}
 a^e b^j c^k\in P
 \Longleftrightarrow
 \begin{cases}
 k\text{ even},\ e\text{ even},\ j>0\text{ even};\quad\text{or}\\
 k\text{ even},\ e\text{ odd},\ j=0;\quad\text{or}\\
 k\text{ odd},\ e\text{ odd},\ j=0\text{ or }j\ge3\text{ odd}.
 \end{cases}
\end{equation}
For completeness, the induction uses only
\[
 a\to1,\qquad b\to\{1,a\},\qquad c\to\{a,b,b^2\}.
\]
Here is the complete case induction.  Put $\epsilon=e\bmod2$.  For a listed
$P$-state with $k$ even, either $(\epsilon,j)=(0,\text{positive even})$
or $(1,0)$; the $a,b,c$ moves land respectively in the complementary
parity classes, and the three $c$ replacements land at
$(1,j),(0,j+1),(0,j+2)$ with $k$ odd.  None is listed.  For a listed state
with $k$ odd, $\epsilon=1$ and $j=0$ or odd at least $3$; the $a$ move
sets $\epsilon=0$, the $b$ moves make $j$ even, and the $c$ moves land at
$(0,j),(1,j+1),(1,j+2)$ with $k$ even.  Again none is listed.

Conversely, every unlisted nonterminal state has the following listed
option.  When $k$ is even:
\begin{center}
\small
\begin{tabular}{ccl}
\toprule
$\epsilon$ & $j$ & move to a $P$-state\\
\midrule
$0$ & $0$ & $c\to a$ (unless $k=0$, when the position is terminal)\\
$0$ & $1$ & $b\to a$\\
$0$ & odd $\ge3$ & $b\to1$\\
$1$ & positive even & $a\to1$\\
$1$ & $1$ & $b\to1$\\
$1$ & odd $\ge3$ & $b\to a$\\
\bottomrule
\end{tabular}
\end{center}
When $k$ is odd, the witnesses are
\begin{center}
\small
\begin{tabular}{ccl@{\qquad}ccl}
\toprule
$\epsilon$&$j$&move&$\epsilon$&$j$&move\\
\midrule
$0$&$0$&$c\to a$&$0$&positive even&$c\to b^2$\\
$0$&odd&$c\to b$&$1$&$1$&$b\to1$\\
$1$&$2$&$c\to a$&$1$&even $\ge4$&$b\to1$\\
\bottomrule
\end{tabular}
\end{center}
The omitted odd-$k$ cases are exactly the listed $P$ cases.  This proves the
mis\`ere recursion for all multiplicities; no exponent cutoff is used.

The outcome law also proves the presentation rather than merely suggesting
it.  It first factors raw monomials through the sixteen forms
$a^\epsilon b^j c^\eta$, with $\epsilon,\eta\in\{0,1\}$ and
$j\in\{0,1,2,3\}$, using
$a^2=c^2=1$ and $b^4=b^2$.  It is invariant under
$b^2\sim ab^3c$.  The congruence generated by this relation pairs the eight
forms in the $b^2$-ideal into four two-element orbits and leaves the eight
forms with $j=0,1$ fixed.  Thus it has exactly twelve candidate classes.
The resulting quotient has presentation
\begin{equation}
 Q_{13}=\langle a,b,c\mid a^2=c^2=1,\ b^4=b^2,\
                 b^2=ab^3c\rangle,
 \qquad P_{13}=\{a,ac,b^2\}.
\end{equation}
It has the twelve normal forms
\[
 1,c,b,bc,b^2,b^2c,b^3,b^3c,a,ac,ab,abc.
\]
Reduction is visible from their distinct $P$-translation rows:
\begin{center}
\small
\begin{tabular}{c@{\;}l@{\qquad}c@{\;}l}
\toprule
$1$&$\{b^2,a,ac\}$&$c$&$\{b^2c,a,ac\}$\\
$b$&$\{b,b^3\}$&$bc$&$\{bc,b^3c\}$\\
$b^2$&$\{1,b^2,abc\}$&$b^2c$&$\{c,b^2c,ab\}$\\
$b^3$&$\{b,b^3,ac\}$&$b^3c$&$\{bc,b^3c,a\}$\\
$a$&$\{1,c,b^3c\}$&$ac$&$\{1,c,b^3\}$\\
$ab$&$\{b^2c,ab\}$&$abc$&$\{b^2,abc\}$\\
\bottomrule
\end{tabular}
\end{center}
Here the row at $x$ is $\{z:xz\in P_{13}\}$.  All rows differ, so the
bipartite monoid is reduced.  Thus the all-multiplicity outcome law proves:

\begin{proposition}\label{prop:grundy13}
The exact mis\`ere quotient generated by $H_1,\ldots,H_{13}$ is the
twelve-element reduced bipartite monoid in the preceding presentation.
\end{proposition}

\subsection{The second extension at heap 18}

The next heaps have values
\[
 x_{14}=b,\quad x_{15}=a,\quad x_{16}=c,\quad x_{17}=b.
\]
Here these assertions are again global, not guesses made inside $Q_{13}$.
With $0$ denoting the terminal game, the split rule gives
\begin{align*}
 \opts(H_{14})&=\{c,ab,a,0\},&
 \opts(H_{15})&=\{c,b,0\},\\
 \opts(H_{16})&=\{ac,b^2,b,a\},&
 \opts(H_{17})&=\{c,ab,a,0\}.
\end{align*}
The reversible-option reductions of Grundy and Smith give respectively
$H_{14}\equiv_-b$, $H_{15}\equiv_-a$, $H_{16}\equiv_-c$, and
$H_{17}\equiv_-b$; these are also the four reductions printed in
\cite[p.~421]{berlekampconwayguy2001}.  Heap $18$ has exact option set
$\{0,b,c,bc\}$ and creates a generator $d=H_{18}$.  By
Lemma~\ref{lem:global-substitution}, every position generated by
$H_1,\ldots,H_{18}$ is globally equivalent to a count state
$a^e b^j c^k d^r$, and every such count state occurs in that universe.

The all-multiplicity recursion extends once more:
\begin{equation}
 \begin{aligned}
 a^e b^j c^k d^r\in P
 \quad\Longleftrightarrow\quad
 &\bigl(r\text{ even and }a^e b^j c^k\text{ satisfies the}\bigr.\\
 &\qquad\bigl.\text{preceding three-generator law}\bigr)\\
 &\text{or }\bigl(r\text{ odd},\ e\text{ even},\ j=0\bigr).
 \end{aligned}
\end{equation}
To check the recursion, use $d\to\{1,b,c,bc\}$.  If $r$ is even and the
three-generator state is $P$, the preceding proof handles the $a,b,c$
moves, while every $d$ move has odd $r$ and either odd $e$ or positive
$j$.  If the three-generator state is $N$, use its displayed winning move.
The only case with no $a,b,c$ component is terminal when $r=0$; for positive
even $r$, the move $d\to1$ reaches the odd-$r$ clause.

If $r$ is odd, a state in the displayed $P$ clause has even $e$ and $j=0$.
An $a$ move makes $e$ odd, a $c$ move makes either $e$ odd or $j$ positive,
and a $d$ move lands, at even $r$, in one of
\[
 (\epsilon,j,k)=(0,0,k),(0,1,k),(0,0,k+1),(0,1,k+1),
\]
all of which fail the three-generator criterion.  Conversely an odd-$r$
$N$-state has a $P$ move in the following exhaustive cases:
\begin{center}
\small
\begin{tabular}{ccl@{\qquad}ccl}
\toprule
$\epsilon$&$j$&move&$\epsilon$&$j$&move\\
\midrule
$1$&$0$&$a\to1$&$0$&$1$&$b\to1$\\
$0$&even $\ge2$&$d\to1$ if $k$ even; $d\to c$ if $k$ odd
 &$0$&odd $\ge3$&$d\to b$ if $k$ even; $d\to bc$ if $k$ odd\\
$1$&$1$&$b\to a$
 &$1$&even $\ge2$&$d\to bc$ if $k$ even; $d\to b$ if $k$ odd\\
$1$&odd $\ge3$&$d\to c$ if $k$ even; $d\to1$ if $k$ odd&&&\\
\bottomrule
\end{tabular}
\end{center}
Each target is one of the three even-$r$ cases already proved.  This
establishes the outcome formula for all $(e,j,k,r)\in\N^4$.

As before, the formula itself identifies the finite presentation.  Before
the last relation there are $32$ forms
$a^\epsilon b^j c^\eta d^\rho$, with
$\epsilon,\eta,\rho\in\{0,1\}$ and $0\le j\le3$, subject to
$a^2=c^2=d^2=1$ and $b^4=b^2$.  The relation
$b^2=ab^3c$ acts freely in the sixteen-element $b^2$-ideal, producing eight
pairs, and leaves the sixteen forms with $j=0,1$ fixed.  Hence exactly
$24$ candidate forms remain, and the outcome law descends to them.  The
quotient is
\begin{equation}
 Q_{18}=Q_{13}\times\langle d\mid d^2=1\rangle,
 \qquad P_{18}=\{a,ac,b^2,d,cd\}.
\end{equation}
Its $24$ normal forms are the twelve above and their products by $d$.
For an explicit reduction certificate, let $R$ be the ordered set of twelve
base forms, put $P_0=\{a,ac,b^2\}$ and $P_1=\{1,c\}$, and define
\[
 A_q=\{z\in R:qz\in P_0\},\qquad
 B_q=\{z\in R:qz\in P_1\}.
\]
The first coordinate $A_q$ is the $Q_{13}$ row already displayed.  The
second coordinates are
\begin{center}
\small
\begin{tabular}{c@{\;}c@{\qquad}c@{\;}c@{\qquad}c@{\;}c}
\toprule
$q$&$B_q$&$q$&$B_q$&$q$&$B_q$\\
\midrule
$1$&$\{1,c\}$&$c$&$\{1,c\}$&$b$&$\varnothing$\\
$bc$&$\varnothing$&$b^2$&$\varnothing$&$b^2c$&$\varnothing$\\
$b^3$&$\varnothing$&$b^3c$&$\varnothing$&$a$&$\{a,ac\}$\\
$ac$&$\{a,ac\}$&$ab$&$\varnothing$&$abc$&$\varnothing$\\
\bottomrule
\end{tabular}
\end{center}
The translation row of $q$ in $Q_{18}$ is $(A_q,B_q)$, while that of $qd$
is $(B_q,A_q)$.  The twelve displayed pairs and their twelve swaps are all
distinct, so every two normal forms have a separating context.  Hence:

\begin{theorem}[Exact Grundy prefix]\label{thm:grundy18}
The exact mis\`ere quotient generated by $H_1,\ldots,H_{18}$ is the
$24$-element reduced bipartite monoid in the preceding presentation.
\end{theorem}

The statements concern entire submonoids generated by the indicated heaps;
no multiplicity bound is present.  They do not decide whether the full
Grundy quotient is finite.  They do show that the apparent period-three
tame pattern fails for an exact, explicit reason at heap $13$, and they give
the first two symbolic extension layers for a future infinite-family or
stabilization argument.

Only the individual-heap reductions displayed above are global.  The
relations $c^2=d^2=1$, $b^4=b^2$, and $b^2=ab^3c$ are conclusions inside
the localized quotient $Q_{18}$; they are not asserted as identities in the
global semigroup of all mis\`ere impartial games.  (The relation $a^2=1$
is global.)

\subsection{A necessary kernel obstruction to finiteness}

The split rule and finiteness alone force a nontrivial recurrent algebraic
component.  This does not decide finiteness, but it excludes stabilization in
any quotient whose kernel is trivial.

\begin{theorem}[Finite Grundy quotients have nontrivial kernel]
\label{thm:grundy-kernel}
Suppose the full mis\`ere quotient $Q$ of Grundy's game is finite.  Then there
is a heap transition $(c,E)$ with $c^2\in E$.  Every two consecutive positive
powers of $c$ are distinct, and the powers of $c$ are eventually periodic
with a period of length at least two.  In particular, the kernel of the finite
commutative monoid $Q$ is nontrivial.
\end{theorem}

\begin{proof}
Color each positive integer $n$ by the heap value $x_n\in Q$.  A standard
distinct-summand form of Schur's theorem gives $i<j$ for which
\[
 x_i=x_j=x_{i+j}=c.
\]
For completeness, it follows directly from finite Ramsey: color the edge
$\{r,s\}$ by $x_{|r-s|}$ and take a monochromatic $K_4$ with vertices
$r_0<r_1<r_2<r_3$.  Splitting the outer difference at $r_1$ and at $r_2$
gives two monochromatic Schur triples; their two summands cannot both be
equal, since that would force $r_1=r_2$.  Choose the unequal one.

Because $i<j$, the position $H_i+H_j$ is a legal option of $H_{i+j}$.
Consequently $c^2\in E$.  Closure of the transition table puts the
$r$-fold transition power
\[
 (c^r,c^{r-1}E)
\]
in the table for every $r\ge1$.  Its option set contains $c^{r+1}$.
Proposition~\ref{prop:noself} therefore gives
$c^r\ne c^{r+1}$ for every $r\ge1$.  Finiteness of $Q$ makes the power
sequence eventually periodic, and the displayed inequalities exclude
period one.

Let $K$ be the kernel, or minimal ideal, of $Q$.  Since $Q$ is a finite
commutative monoid, $K$ is a group; write $z$ for its identity.  Surjectivity
of the quotient map supplies a transition with value $z$.  Multiplying it by
$(c,E)$ gives a kernel-valued transition of value $k=zc$ whose option set
contains
\[
 zc^2=(zc)^2=k^2.
\]
Again no-self gives $k\ne k^2$.  A one-element kernel would have $k=z=k^2$,
a contradiction.
\end{proof}

The conclusion is structural rather than contradictory: period two already
occurs in tame quotients.  To prove the full Grundy quotient infinite, one
still needs either unbounded kernel rank or infinitely many distinct
translation rows, not merely one nonidentity kernel element.

\section{Formal verification and open boundary}\label{sec:boundary}

The results divide the original problem into sharper pieces.
\begin{enumerate}
\item Every tame finite quotient $\mathcal T_n$ has a uniform finite-code
      realization.  For the $|P|=2$ classification, the natural realization
      question remains for the family obtained by tame extension of
      $\mathcal R_8$.
\item Fix a finite $(Q,P)$, a valid table $T$, and a finite code $\Gamma$.
      At stage $n$, the preceding values determine $E_n$, and determinism
      supplies at most one possible $x_n$.  Exact realization through $T$ is
      therefore precisely totality of this partial recurrence plus generation
      of $Q$.  For no-split codes this is finite-state.  With splitting, the
      first unresolved problem is deciding totality of the unbounded
      convolution recurrence for a fixed $T$ and $\Gamma$.  Theorem~
      \ref{thm:decoder-hardness} proves $\Pi^0_1$-completeness when the decoder
      is arbitrary, so a positive algorithm must exploit validity; the theorem
      does not transfer hardness to valid tables.
\item Even a decision theorem for every fixed split code would leave a second
      problem: synthesis ranges over codes of unbounded length.  A decidable
      monoid-level characterization needs a computable length bound or pumping
      normalization, unless natural realization is instead proved
      undecidable.  There are only finitely many tables on fixed $Q$, so table
      enumeration introduces no third gap.
\item For Grundy's game, $Q_{13}$ and $Q_{18}$ are exact.  A solution of the
      full problem must either continue these presentations uniformly and
      prove stabilization, or produce an infinite family of distinct
      translation rows with explicit separating contexts.  In the finite
      branch, Theorem~\ref{thm:grundy-kernel} now additionally requires a
      nontrivial kernel and nonconstant eventual power cycle.
\end{enumerate}

The import-only entry point
\texttt{Ogdoad.Papers.}\allowbreak\texttt{MisereNaturalRealization}
exposes the Lean development.  It checks transition determinism, meximal
inclusion, $P$-reaching contexts, the no-self-option theorem and all three
octal digit exclusions.  It also checks the generic rank/parity proof that a
surjective multiplicative value map has exactly the quotient fibers claimed,
the complete-option-record periodicity bound and finite infinite-tail
certificate, and the abstract prefix/pad/unary-tail construction underlying
Theorem~\ref{thm:normalform}.  It also checks the inert-pad obstruction,
the coefficientwise unary-language dictionary and failure language, the
distinct-summand finite-color bridge (via Mathlib's Hindman theorem), and the
algebraic implication from a square option to a nonconstant eventual power
period.  The grammar-hardness reduction and kernel deduction are paper-level,
as are the transport from typed prefix
heaps to their numerical rank-order enumeration, and passage from ordered
heap lists to multisets.  The tame-family strategy and the Grundy prefix
presentations likewise remain paper proofs; their finite monoid models and
exponent recursions are the next natural formalization targets.

\section{Conclusion}

Natural realization is not an additional finite-table axiom.  It is a
totality problem for a quotient-valued heap trace.  Transition-table
determinism makes that trace canonical: no-split codes become finite-state,
while split codes expose exactly the convolution that prevents the same
automatic periodicity argument.  The unary-language reduction shows that this
is not a superficial failure of the proof: without valid-table structure, the
fixed-code totality problem is already $\Pi^0_1$-complete.

The finite-exception normal form shows that a uniform infinite tail alone is
not the difficulty: every finite quotient admits such a realization.  The
octal constraint couples each intended finite-prefix edge to infinitely many
translated edges, and the no-self inequalities give the first absolute
rejection test for that coupling.  The inert-pad theorem rules out the most
direct compiler, but leaves active bridges open.  What remains is to decide
totality under the full valid-table axioms and then control unbounded code
synthesis.

Within this framework the tame family is complete.  One explicit sequence of
finite codes realizes every $\mathcal T_n$, with an outcome proof valid for
all positions.  Grundy's game exhibits the complementary phenomenon: its
tame prefix breaks at heap $13$ and extends again at heap $18$, but each
finite stage remains exactly computable as a reduced all-multiplicity
quotient.  Any finite limit must have a nontrivial kernel, but that condition
is compatible with period two.  A complete solution must now control the
extension process uniformly---either by eventual stabilization or by explicit
separating contexts for infinitely many classes.

\bibliographystyle{plain}
\bibliography{misere_natural_realization}

\end{document}
